\documentclass[11pt]{article}
\usepackage[T1]{fontenc}
\usepackage[textwidth=250pt]{geometry}
\usepackage{microtype}
\pagestyle{empty}
\begin{document}
The first sum is $(1+2)=3$ and the second one is $(4+5)=9$, while the
product $2(3+4)=14$ is checked by hand. A set such as $[0,1]$ is closed, and
the interval $(0,1)$ is open; together they give $(2+3)(4+5)=45$ exactly.

We compute $7+8=15$ and then $15-6=9$ before we write the answer down. The
next step gives $9+10=19$, and adding $1$ to it yields $20$, which is the
total number of points available on this particular problem set.

Numbers such as $2024$ and $365$ appear in the text, and a longer expression
like $(1+2+3+4+5+6)=21$ sits in the middle of a line that must stretch or
shrink a little so that both of its margins stay perfectly straight.
\end{document}
